% !TEX root = ../main.tex
% ==============================================================================
% ANEXO OBLIGATORIO: BORRADOR DE CARTA DE APOYO SIMULADA
% Responsable: Anais (Vinculación Institucional)
% ==============================================================================

\section{Anexo: Borrador de Carta de Apoyo Institucional Simulada}
\label{sec:carta_apoyo}

\noindent
\textit{\small \textbf{Nota de contextualización académica:} El siguiente documento simula la carta de apoyo exigida en la Fase 3, asumiendo una alianza con \textbf{MCI Electronics}. Su rol como red de apoyo externa consiste en brindar asesoría técnica para la configuración de los microcontroladores y otorgar un 15\% de descuento en el hardware, fortaleciendo la viabilidad del proyecto.}

\vspace{0.8em}

\vspace{1cm}

\begin{flushright}
\textbf{MCI ELECTRONICS}\\
Soluciones Tecnológicas Integrales\\
Valparaíso, 4 de noviembre de 2026\\
\vspace{0.5cm}
\textbf{REF:} Patrocinio y colaboración técnica al proyecto ``Campus Zero''.
\end{flushright}

\vspace{1cm}

\noindent
\textbf{A: Equipo de Proyecto ``Campus Zero''}\\
Asignatura INF276 - Ingeniería, Informática y Sociedad\\
Universidad Técnica Federico Santa María

\vspace{1cm}

\noindent Estimados estudiantes:

\vspace{0.5cm}

\noindent Por medio de la presente, desde MCI Electronics queremos expresar nuestro genuino interés y respaldo a su iniciativa \textbf{``Campus Zero: Control Automatizado de Iluminación y Climatización de Salas''}. Como empresa proveedora de insumos y soluciones de base tecnológica, valoramos profundamente que los futuros ingenieros de la UTFSM busquen resolver problemas reales de eficiencia energética y sustentabilidad utilizando hardware libre.

\vspace{0.3cm}

\noindent Sabemos que la implementación de un piloto requiere sortear diversas barreras técnicas y de presupuesto. Por ello, para fomentar este tipo de innovaciones en el ecosistema estudiantil, nos comprometemos a apoyar activamente su proyecto en caso de adjudicarse los fondos, mediante las siguientes acciones:

\begin{itemize}
    \item \textbf{Asesoría Técnica Especializada:} Facilitaremos una sesión de trabajo telemática con uno de nuestros ingenieros de soporte. El objetivo será validar la arquitectura electrónica que diseñaron y asesorar en la correcta configuración y buenas prácticas de los microcontroladores \textbf{ESP32 Wi-Fi}, previniendo fallos en terreno.
    \item \textbf{Descuento Estudiantil de Apoyo:} Aplicaremos un descuento preferencial del 15\% sobre el precio de lista para la compra de las placas ESP32, relés, fuentes y sensores detallados en su presupuesto, ayudándoles a resguardar su fondo de contingencia.
\end{itemize}

\vspace{0.3cm}

\noindent Creemos plenamente en la capacidad de su equipo y en el impacto medioambiental positivo que Campus Zero puede lograr en su universidad. Quedamos a su entera disposición para comenzar el trabajo conjunto.

\vspace{1cm}

\noindent Atentamente,

\vspace{2cm}

\begin{center}
\rule{7cm}{0.5pt}\\
\textbf{[Borrador Simulado - Sin Firma]}\\
\vspace{0.2cm}
Área de Vinculación y Proyectos Educativos\\
MCI Electronics Chile
\end{center}
